\section{Deformations of the cut data}
\label{sec:deformation}

The cut data in Proposition~\ref{geo:construction} have two consequences
for representations near $\rho$: complementary tangent spaces for the
paired wall restrictions, and a cube-zero ideal generated by the end
holonomies minus the identity.  The first comes from the graph homotopy
type of $G$ and the longitude--meridian interchange; the second uses the
triple commutator relations retained in the component $X_Q$.

We use $\Gamma$ and $\Gamma_h$ for the graph and its $h$-subrose.
The coefficient system $\mathfrak l_\rho$ is the pullback of the adjoint
system of \eqref{geo:matrices}, extended by $\rho(q_i)=1$.  It is trivial
on every wall, with matching reference frames on the two copies induced
by the cut point.  The trace form
$(u,v)\mapsto\operatorname{tr}(uv)$ identifies this system with its dual.
Throughout, cohomology of $X$ or $J_i$ includes all components: $X$ need
not be connected, and a wall may have closed components in addition to
its one component with torus boundary.

\subsection{The paired restriction maps}

\begin{lemma}[Paired Mayer--Vietoris maps]
\label{def:paired-mv}
Let
\[
 W_h=\operatorname{im}\bigl(H^1(\Gamma_h;\mathfrak l_\rho)
                 \longrightarrow H^1(X;\mathfrak l_\rho)\bigr).
\]
Here the map uses the retraction of the cut-open graph onto its
\(h\)-rose.  For either degree, write \(r_i^\pm\) for restriction to
the two copies of \(J_i\).  The map
\begin{equation}
 \Delta(u,v)=
  \bigl(r_i^+u-r_i^-v,\ r_i^-u-r_i^+v\bigr)_{i=1}^d
 \label{def:delta}
\end{equation}
is onto in degree one, with kernel \(W_h\oplus W_h\), and is an
isomorphism in degree two.  Consequently, if
\(H^1(X;\mathfrak l_\rho)=W_h\oplus U\), then
\[
 \Delta:U\oplus U\xrightarrow{\ \cong\ }
             \bigoplus_i\bigl(H^1(J_i;\mathfrak l)
                          \oplus H^1(J_i;\mathfrak l)\bigr).
\]
\end{lemma}

\begin{proof}
Take the two-sheeted cover of \(G\) classified by the homomorphism
sending every \(q_i\) to the nontrivial element of \(\mathbb Z/2\)
and every \(h\)-loop to zero.  The lift of \(p\) is a homotopy
equivalence with the corresponding graph cover.  Thus the cover has
zero cohomology in degrees at least two, with the lifted local
system as well as with constant coefficients.  After adjoining
small product collars, its decomposition has two vertex spaces,
both copies of \(X\), and two edge spaces per label, both copies of
\(J_i\).  The pairings go from one vertex copy to the opposite side
of the other, giving exactly \eqref{def:delta}.  The relevant part of
Mayer--Vietoris is
\[
 H^1(\widetilde G)\longrightarrow H^1(X)^{\oplus2}
 \xrightarrow{\Delta}\bigoplus_i H^1(J_i)^{\oplus2}
 \longrightarrow H^2(\widetilde G)=0
\]
and
\[
 0=H^2(\widetilde G)\longrightarrow H^2(X)^{\oplus2}
 \xrightarrow{\Delta}\bigoplus_i H^2(J_i)^{\oplus2}
 \longrightarrow H^3(\widetilde G)=0 .
\]
The omitted coefficients in these two displays are the specified
local systems.

The image in the middle degree-one group is precisely
\(W_h\oplus W_h\).  In one direction, a class restricted from the
graph cover factors on each vertex through its cut-open graph,
and hence through the associated \(h\)-rose.  Conversely, choose
cellular one-cochains representing any two classes on the two
\(h\)-roses.  Extend them to the graph cover by arbitrary values,
for example zero, on the remaining edges.  There are no two-cells,
so the extensions are cocycles.  Their restrictions give the
chosen pair of classes on \(X\).  This argument uses direct sums
over all components of \(X\), and imposes no connectedness
assumption on it.
\end{proof}

\begin{lemma}[Slopes and internal directions]
\label{def:slopes}
The ends of each cut can be designated \(+\) and \(-\), and torus
coordinates \((\theta_i,\eta_i)\) of period one can be chosen, so that
the image of
\[
 H^1(J_i;\mathbb C)\longrightarrow H^1(\partial J_i;\mathbb C)
\]
is the \(\theta_i\)-axis.  In \(Q\), the \(\eta_i\)-circle bounds on
the positive side and the \(\theta_i\)-circle bounds on the negative
side.  Put
\[
 n_i=\dim_{\mathbb C}H^1(J_i,\partial J_i;\mathfrak l),
 \qquad D=\sum_i n_i,
 \qquad
 \mathcal Y=\bigoplus_i(\mathbb C^{n_i}\oplus\mathbb C^{n_i}).
\]
Then \(\dim U=3d+D\).  The joint first-slope restriction
\(U\to\mathfrak l^d\) is onto.  Its kernel \(U_0\) has dimension \(D\).
The two maps from \(U_0\) to \(\mathcal Y\) given by the paired
restrictions and by their interchange are injective and have
complementary images.
\end{lemma}

\begin{proof}
Poincare--Lefschetz duality \cite[Theorem~3.43]{Hatcher2002} and the exact sequence of the pair imply
that the restriction image in the cohomology of the boundary of
an oriented compact three-manifold is its own annihilator for
the intersection form.  For completeness, the connecting map
\(H^1(\partial J_i)\to H^2(J_i,\partial J_i)\) is dual, up to sign,
to the restriction map.  Exactness therefore identifies the
kernel of that connecting map, which is the restriction image,
with the annihilator of that image.  On a single torus a subspace
equal to its annihilator is a line.  Closed components of \(J_i\)
make no contribution to the boundary map.  Denote this line by
\(L_i\).

Tensoring with \(\mathfrak l\) gives the corresponding assertion
for the trivial wall system.  A restriction from the positive or
negative copy of \(X\) has zero evaluation on the slope which
bounds in \(Q\), because the local system is trivial on that
torus.  These two constraints are the distinct coordinate axes
interchanged by the Hopf gluing.  If \(L_i\) were neither axis,
both restrictions would have zero torus value.  This contradicts
surjectivity in Lemma~\ref{def:paired-mv}, since
\(H^1(J_i;\mathfrak l)\) has a nonzero boundary restriction.
Thus \(L_i\) is one axis; designate the contributing end positive
and call this the first slope.

The map \(H^0(J_i;\mathfrak l)\to
H^0(\partial J_i;\mathfrak l)\) is onto: constants on the unique
component with boundary suffice.  The exact sequence consequently
identifies
\begin{equation}
 H^1(J_i,\partial J_i;\mathfrak l)
 \ \cong\
 \ker\bigl(H^1(J_i;\mathfrak l)
                  \longrightarrow H^1(\partial J_i;\mathfrak l)\bigr).
 \label{def:relative-identification}
\end{equation}
These are the \(n_i\) internal directions, including all degree-one
cohomology of every closed component.  They pair perfectly with
\(H^2(J_i;\mathfrak l)\) by integration and the trace form.
In particular \(\dim H^1(J_i;\mathfrak l)=3+n_i\).

The isomorphism in Lemma~\ref{def:paired-mv} now gives
\(\dim U=3d+D\).  Its first-slope coordinates are the first-slope
values of \(u\), and the negatives of those of \(v\), because all
opposite-side values are zero.  Surjectivity implies independent
surjectivity onto the \(3d\) first-slope coordinates.  On their
kernel the same isomorphism becomes
\[
 U_0\oplus U_0\longrightarrow\mathcal Y,\qquad
 (u,v)\longmapsto R(u)-\sigma R(v),
\]
where \(R\) is paired restriction and \(\sigma\) interchanges the
two slots for each \(i\).  This proves the last assertions.
\end{proof}

\subsection{Finite algebraic equation charts}

From now on orient \(J_i\), including each of its closed components,
as its positive copy in the oriented boundary of \(X\).  Its
negative copy then has the opposite orientation.

\begin{definition}[Scalar rings and formal coefficients]
\label{def:coefficient-convention}
An algebraic parameter chart will mean a local \(\mathbb C\)-algebra
essentially of finite type and etale over an affine coordinate
space at the designated origin.  Its completion is
\(B=\mathbb C[[v_1,\ldots,v_N]]\), with maximal ideal \(\mathfrak m\).
Finite collections of algebraic germs obtained by an implicit
equation with invertible Jacobian can always be placed in one such
chart, after replacing it by a common etale neighborhood.

For a smooth manifold \(Z\), smooth forms with formal coefficients
mean
\[
 \Omega^j(Z)\widehat\otimes B
   =\varprojlim_k\bigl(\Omega^j(Z)\otimes_{\mathbb C}
                                     B/\mathfrak m^k\bigr).
\]
For quotients of \(B\), use the analogous inverse limit.  The
extension of a scalar ideal \(I\) to this space always means its
closed extension.  Differential operations in the \(Z\)-variables
preserve that extension; multiplication respects its powers.
Integration over a compact manifold is defined at each finite
parameter quotient and then by inverse limit.
\end{definition}

The tangent calculation alone does not retain the equations that obstruct
a deformation.  We now choose finite algebraic transition charts and
keep their entire defect ideals, including their nonreduced structure.
The construction removes equations with independent linear terms but
retains every higher-order defect.  A nonzero functional on its first ideal
layer $\mathfrak r_Z/\mathfrak m\mathfrak r_Z$ detects an obstruction to
a flat lift, even when that lift may leave the chosen tangent slice.  This is
the property used later to identify the full derivative ideal in
Proposition~\ref{conn:derivative-ideal}.  Small extensions and their
second-cohomology obstructions are part of the deformation framework of
\cite[Proposition~2.6(1)]{GM1988}; the explicit construction below keeps
the triangle-defect ideal needed here.

\begin{lemma}[Equation charts and their minimal defect layer]
\label{def:equation-chart}
Let \(Z\) be a compact smooth manifold, possibly disconnected and
with boundary, and let \(\rho_Z\) be a fixed
\(\operatorname{SL}_2(\mathbb C)\)-local system.  Choose a subspace
\(T\subset H^1(Z;\mathfrak l_{\rho_Z})\).
There is an algebraic family of edge matrices on a finite good
cover, with parameters \(T\), having these properties:
\begin{enumerate}
\item Its reference value is the chosen local system, its tangent
      is the specified inclusion of \(T\), and its triangle-defect
      ideal \(\mathfrak r_Z\) is contained in \(\mathfrak m^2\).
\item Modulo \(\mathfrak r_Z\), its matrices are the transitions of
      an exact flat family.
\item For a nonzero functional
      \(\lambda:\mathfrak r_Z/\mathfrak m\mathfrak r_Z\to\mathbb C\),
      let \(\mathfrak s_\lambda\subset\mathfrak r_Z\) be the inverse image
      of \(\ker\lambda\).  The specified flat family over
      \(B/\mathfrak r_Z\), where \(B=\mathbb C[[T]]\), has no flat lift
      over \(B/\mathfrak s_\lambda\).  This excludes compatible
      general-linear lifts as well, and does not require a putative
      lift to remain in the parameter slice \(T\).
\end{enumerate}
\end{lemma}

\begin{proof}
Use a finite good cover, with contractible connected nonempty
intersections, obtained from a finite smooth triangulation and
boundary collars.  Components are covered independently.
Fix reference parallel frames.  For each oriented edge of the
nerve, vary a matrix \(g_{\alpha\beta}\) in an algebraic group
coordinate patch, impose
\(g_{\beta\alpha}=g_{\alpha\beta}^{-1}\), and set
\(g_{\alpha\alpha}=1\).  For an oriented triangle use the
traceless part of
\[
 g_{\alpha\beta}g_{\beta\gamma}g_{\gamma\alpha}-1
\]
as its three equation coordinates.  These coordinates generate
the full matrix defect ideal locally.  Indeed, write the triangle
product as \((1+x)1+T\), where \(T\) is traceless and \(x(0)=0\).
Its determinant equation reads
\(x(2+x)=-\det T\).  The factor \(2+x\) is a unit, so \(x\)
belongs to the ideal generated by the entries of \(T\)
(in fact to its square).

In reference frames, the linearized defect map is the cochain
differential
\[
 d_\rho:C^1(Z;\mathfrak l_{\rho_Z})
                    \longrightarrow C^2(Z;\mathfrak l_{\rho_Z})
\]
on the good-cover nerve.  Choose decompositions
\[
 C^1=\operatorname{im}d_\rho^0\oplus H\oplus L,\qquad
 \ker d_\rho^1=\operatorname{im}d_\rho^0\oplus H,
\]
and choose a projection
\(\pi:C^2\to\operatorname{im}d_\rho^1\) which is the identity on
that image.  Set the first summand of the coordinate variables
equal to zero.  Solve \(\pi(\text{defect})=0\) for the
\(L\)-variables.  Its derivative in those variables is the
isomorphism
\(d_\rho^1:L\to\operatorname{im}d_\rho^1\).
The algebraic implicit function theorem gives the solution in
an etale neighborhood.  Its remaining coordinates are \(H\),
and its defect has no linear term.  Finally restrict the
\(H\)-parameters to the chosen linear subspace \(T\).
The identity \(\pi(\text{defect})=0\) survives this restriction.
Good-cover cohomology agrees with local-system cohomology:
the cover and every intersection are acyclic for these locally
constant coefficients.  Thus these are the asserted tangent
coordinates.

Write \(B=\mathbb C[[T]]\), \(\mathfrak r=\mathfrak r_Z\), and
suppose \(\lambda\ne0\).  Its inverse-image kernel is the ideal
\[
 \mathfrak s
 =\{b\in\mathfrak r:\lambda(b\bmod\mathfrak m\mathfrak r)=0\}.
\]
It is an ideal because multiplication acts on
\(\mathfrak r/\mathfrak m\mathfrak r\) by its constant term.
There is an exact sequence
\begin{equation}
 0\longrightarrow I=\mathfrak r/\mathfrak s
 \longrightarrow B/\mathfrak s
 \longrightarrow B/\mathfrak r\longrightarrow0,
 \qquad
 \dim_{\mathbb C}I=1,\quad \mathfrak m I=I^2=0.
 \label{def:small-extension}
\end{equation}
Here \(\mathfrak r^2\subset\mathfrak m\mathfrak r\) proves the last
equality.  Evaluating the triangle defects in \(I\) gives a
reference two-cochain \(c\) with \(\pi c=0\).
It is nonzero, since the entries of the defects generate
\(\mathfrak r\) and hence span its quotient by
\(\mathfrak m\mathfrak r\).

A flat lift changes edge matrices by \(I\)-valued corrections.
Since \(\mathfrak mI=0\), their effect on triangle products is
the reference differential, so flatness would give
\(c+d_\rho u=0\).  Applying \(\pi\) gives \(d_\rho u=0\), and
then \(c=0\), a contradiction.  If the corrections are
general-linear, project them to their traceless parts.
Conjugation preserves the splitting into scalar and traceless
matrices, so this projection commutes with \(d_\rho\);
the same contradiction follows.  If the putative flat lift is
specified only up to an isomorphism of its reduction, choose
parallel frames on the contractible patches with initial
values lifting the given reduced frames.  Their transition
matrices then lift exactly the original matrices, which is the
case just proved.
\end{proof}

\begin{remark}[Based wall charts]
\label{def:based-wall}
At the trivial wall representation the preceding construction can
be made in based tree frames: choose a maximal tree in each
component of the one-skeleton of the nerve and set its edge
matrices equal to \(1\).  This removes exactly the degree-one
coboundary directions, without dividing by constant conjugation.
Any based flat representation sufficiently near the reference
one, including one over a complete local quotient, is represented
by its edge matrices in these tree frames.  It lies in the implicit
chart: take its free \(H^1\)-coordinates and use uniqueness in the
solved \(L\)-coordinates.  The same assertion follows order by
order over an arbitrary complete quotient because the initial
Jacobian determinant remains a unit.
\end{remark}

\begin{proposition}[Wall and vertex coordinates]
\label{def:restriction-charts}
There are wall charts with parameters
\[
 (a_i,y_i)\in\mathbb C^3\oplus\mathbb C^{n_i},
\]
vertex parameters
\((a,t)\in\mathbb C^{3d}\oplus\mathbb C^D\), a vertex defect ideal
\(\mathfrak r\subset\mathbb C[[a,t]]\), and algebraic maps
\begin{equation}
 \iota_i^+(a,t)=(a_i,Y_i^+(a,t)),\qquad
 \iota_i^-(a,t)=(0,Y_i^-(a,t)),
 \label{def:restriction-maps}
\end{equation}
with the following properties.  The maps give the actual based
restrictions of the flat vertex family modulo \(\mathfrak r\).
The \(a_i\) are algebraic group coordinates for the first-slope
holonomy, centered at the identity.  The two maps
\[
 t\longmapsto Y(0,t),\qquad
 t\longmapsto \sigma Y(0,t),
 \quad Y=(Y_i^+,Y_i^-)_{i=1}^d,
\]
have injective and complementary tangent images in \(\mathcal Y\).
\end{proposition}

\begin{proof}
Apply Lemma~\ref{def:equation-chart} to each \(J_i\) at its trivial
reference representation, taking the full \(H^1\), in the tree
frames of Remark~\ref{def:based-wall}.  A fixed edge word for the
first torus generator gives an algebraic matrix function even
when the transition cochain has nonzero defects.  By
Lemma~\ref{def:slopes} its three group coordinates have independent
linear terms.  Replace three of the \(H^1\)-coordinates by these
coordinates, preserving complementary \(y_i\)-coordinates.
The algebraic inverse function theorem permits this change.

Apply Lemma~\ref{def:equation-chart} to \(X\), retaining \(U\).
Restriction of a flat cochain to a wall, followed by tree
normalization, is computed by finitely many based path words.
One can arrange compatible refinements of the covers, or
subdivide each of these finitely many paths into patches; in
either description all resulting matrix expressions are finite
products of transitions and their inverses.  Taking the free
coordinates of the wall chart therefore gives algebraic lifts of
the restriction parameters.  Over the quotient by \(\mathfrak r\)
the solved coordinates reproduce the same cochain by the
uniqueness just proved.

The joint positive first-slope values on \(U\) have surjective
derivative.  Use their actual group coordinates as \(a\), retaining
complementary coordinates \(t\).  On the negative side the first
slope bounds in \(Q\), so its holonomy is exactly \(1\) modulo
\(\mathfrak r\); its coordinate lift may thus be set equal to
zero.  This gives \eqref{def:restriction-maps}.  Basing comparisons
can also be lifted algebraically: they are products along finitely
many paths and changes between chosen local frames.  At reference
all wall holonomies are \(1\), so a varying basing frame makes
no first-order change to their conjugated holonomies.  The
tangent maps are consequently the ordinary restriction maps.
Lemma~\ref{def:slopes} gives their asserted complementarity.
\end{proof}

Let \(H_{i,1}(a_i)\) and \(H_{i,2}(a_i,y_i)\) be the two slope
matrices computed by the chosen words.  Write \(\mathfrak d_i\)
for the wall triangle-defect ideal and put
\begin{equation}
 U_i=-\log H_{i,1},\qquad V_i=-\log H_{i,2},\qquad
 \mathfrak k_i=
  \bigl(\mathfrak d_i,\ H_{i,1}-1,\ H_{i,2}-1\bigr).
 \label{def:wall-ideal}
\end{equation}
Matrix entries generate the ideals in this notation.
The logarithms are formal matrix logarithms; no algebraicity
of their full series is asserted.  They are traceless, because
the matrices have determinant one and reference value \(1\).
The second-slope restriction has zero derivative, so
\[
 V_i\in\mathfrak m_i^2,\qquad \mathfrak d_i\subset\mathfrak m_i^2.
\]
The first slope depends exactly on \(a_i\), not on \(y_i\).
The convention \(U_i=-\log H_{i,1}\) is suited to the connection
convention \(d+A\), whose constant-circle holonomy is
\(\exp(-A(\partial_\theta))\).

\subsection{The nilpotent ideal of end holonomies}

We now use the full finite word list.  In $\pi_1(X_Q)$ the geometric
construction gives
\begin{equation}
 [z_1^{c_1},[z_2^{c_2},z_3^{c_3}]]=1,
 \qquad c_1,c_2,c_3\in\mathcal H,
 \label{def:relations}
\end{equation}
where the three end labels are arbitrary and may repeat.  Their
independent conjugators turn the leading commutator terms into every
scalar cubic product.  The following calculation explains why words of
length at most five suffice.

\begin{lemma}[A finite adjoint spanning set]\label{geo:spanning}
For the word set $\mathcal H$ fixed in Section~\ref{sec:geometry} and
the representation in \eqref{geo:matrices}, one has
\[
 \operatorname{span}_{\mathbb C}
 \{\operatorname{Ad}\rho(c):c\in\mathcal H\}
       =\operatorname{End}_{\mathbb C}(\mathfrak l).
\]
\end{lemma}

\begin{proof}
Use the ordered basis
$E=\left(\begin{smallmatrix}0&1\\0&0\end{smallmatrix}\right)$,
$H=\left(\begin{smallmatrix}1&0\\0&-1\end{smallmatrix}\right)$,
$F=\left(\begin{smallmatrix}0&0\\1&0\end{smallmatrix}\right)$.
The operator $A=\operatorname{Ad}\rho(h_0)$ has distinct eigenvalues
$4,1,1/4$.  Its three eigenprojections $P_E,P_H,P_F$ are polynomials of
degree at most two in $A$.  If $B_\pm=\operatorname{Ad}\rho(h_\pm)$, then
\[
 \begin{array}{lll}
 B_+E=E,&B_+H=H-2E,&B_+F=F+H-E,\\
 B_-E=E-H-F,&B_-H=H+2F,&B_-F=F.
 \end{array}
\]
For every two distinct basis lines, one of
$P_iB_+P_j,P_iB_-P_j$ is a nonzero multiple of the corresponding
matrix unit.  These operators are linear combinations of
$A^aB_\pm A^b$ with $0\leq a,b\leq2$, hence of adjoints of words of
length at most five.  The diagonal matrix units are the $P_i$ themselves.
Thus the adjoints indexed by $\mathcal H$ span all nine matrix units.
\end{proof}

\begin{proposition}[Cubic nilpotence]
\label{def:cubic-nilpotence}
In \(B/\mathfrak r\), where \(B=\mathbb C[[a,t]]\), let
\(\mathfrak n\) be generated by all matrix entries of the end
holonomies minus \(1\).  Then
\begin{equation}
 \mathfrak n^3=0.
 \label{def:n-cube}
\end{equation}
In particular, for either restriction map in
\eqref{def:restriction-maps},
\begin{equation}
 (\iota_i^\pm)^*(\mathfrak k_i)^3\subset\mathfrak r .
 \label{def:k-cube}
\end{equation}
The same conclusion holds for the cube of the sum of all these
pulled-back ideals.
\end{proposition}

\begin{proof}
All computations initially take place in the Noetherian local
ring \(B/\mathfrak r\).  Let \(Z=z-1\) for an end matrix.  The
determinant equation gives
\(\operatorname{tr}Z=-\det Z\in\mathfrak n^2\).
Thus one may replace \(Z\) by its traceless part when computing
triple products modulo \(\mathfrak m\mathfrak n^3\).
Indeed the discarded products lie in
\(\mathfrak n^4\subset\mathfrak m\mathfrak n^3\).

The formal group-commutator expansion is
\[
 [1+A,[1+B,1+C]]
     =1+[A,[B,C]]+\text{terms of total degree at least four}.
\]
It follows by expanding each inverse as
\((1+A)^{-1}=1-A+A^2-\cdots\); the inner commutator starts with
\(1+[B,C]\), and its commutator with \(1+A\) first contributes
\([A,[B,C]]\).  In \eqref{def:relations}, replacing a conjugating
matrix by its reference value changes the displayed cubic term
only by \(\mathfrak m\mathfrak n^3\).  Every relation therefore
implies, in \(\mathfrak n^3/\mathfrak m\mathfrak n^3\),
\[
 [\operatorname{Ad}\rho(c_1)Z_1,
   [\operatorname{Ad}\rho(c_2)Z_2,
    \operatorname{Ad}\rho(c_3)Z_3]]=0,
\]
with traceless \(Z_j\).
The independent choice of all three conjugators, followed by
linear combination and Lemma~\ref{geo:spanning}, replaces the three
adjoint maps by arbitrary endomorphisms of \(\mathfrak l\).

For any three scalar linear functionals \(\ell_1,\ell_2,\ell_3\)
on \(\mathfrak l\), choose these endomorphisms as
\[
 x\longmapsto\ell_1(x)E,\qquad
 x\longmapsto\ell_2(x)H,\qquad
 x\longmapsto\ell_3(x)F,
\]
where
\[
 E=\begin{pmatrix}0&1\\0&0\end{pmatrix},\quad
 H=\begin{pmatrix}1&0\\0&-1\end{pmatrix},\quad
 F=\begin{pmatrix}0&0\\1&0\end{pmatrix}.
\]
Since \([E,[H,F]]=-2H\ne0\), every product
\(\ell_1(Z_1)\ell_2(Z_2)\ell_3(Z_3)\) vanishes in this quotient.
The relation list permits repeated ends, so this includes all
generators of \(\mathfrak n^3/\mathfrak m\mathfrak n^3\).
Hence \(\mathfrak n^3=\mathfrak m\mathfrak n^3\); Nakayama's lemma,
applied to the finitely generated ideal \(\mathfrak n^3\), proves
\eqref{def:n-cube}.

Under either restriction, the wall defects vanish modulo
\(\mathfrak r\).  Each slope holonomy is either \(1\) or the
holonomy of its end longitude, with a possible inverse and a
change of basing.  Entries of its deviation therefore belong to
\(\mathfrak n\).  This remains true for inverses, since
\(z^{-1}-1=-z^{-1}(z-1)\), and for conjugation.  The images of all
the \(\mathfrak k_i\) are consequently contained in
\(\mathfrak n\), proving the final statements.
\end{proof}
